\chapter{Legislación}
Apuntes sobre legislación concreta. Hay muchas leyes que salen una y otra vez. No es temario \textit{ex profeso} pero si merece la pena una pasada por ellas.

\section{Sobre las leyes}
Las leyes tienen su orígen en la soberanía popular a través de las cortes. (Bien para su aprobación, bien para su confirmación) y las firma el Rey.

Los principios descritos pueden no aplicarse a todos los casos. Pues tenemos en España una tradición legislativa muy amplia que ha cambiado con el tiempo.

\subsection{Prioridad}
    El escalafón de las leyes es el que sigue. Teniendo prioridad las leyes de categoría superior sobre la inferior y las nuevas sobre las viejas. La categoría viene dada, fundamentalmente, por la dificultad para aprobar o derogar una ley.
    \begin{description}
        \item[Constitución Española] Base del derecho en españa.
        \item[Derecho de la Unión europea] España aplica los reglamentos por ser estado miembro sin necesidad de transposición. Los estdos miembros deben aplicar las directivas, pero tienen libertad para hacerlo mediante leyes propias.
        \item[Tratados internacionales] Aprobados en las cortes, no sólo por el hecho de existir.
        \item[Ley Orgánica (LO) ] Desarrolla derechos fundamentales y libertades públicas, el régimen electoral y otros aspectos previstos en la constitución. Su modificación exige mayoría absoluta.
        \item[Ordinaria (L)] Emana de las cortes generales y se aprueba por mayoría simple.
        \item[Especial (LE)] Casos como los estatutos de autonomía, que se aprueban como ley orgánica.
        \item[Decreto-ley (DLx)] Emana del gobierno en lugar del parlamento. Se aprueba en casos de necesidad y posteriormente se someten al contol del parlamento. No pueden sustituir a leyes orgánicas.
        \item[Decreto legislativo (DLeg)] Emana del gobierno por delegación de las cortes, para dictar textos articulados o refundidos. No pueden regular materias que corresponden a leyes orgánicas
        \item[Ley autonómica(LA)] Emana de los parlamentos autonómicos. Solo puede reglar las competencias atribuidas a la Comunidad Autónoma.
        \item[Real decreto (RD)] (No tiene rango de ley) Reglamento emitido por consejo de ministros  presidente.
        \item[Orden ministerial (OM)] (No tiene rango de ley) Reglamento dictado por un ministro.
        \item[Reglamentos] Existen multitud de tipos de normas en función de dónde emanan.
    \end{description}

\subsection{Forma}
Las normas, siempre deben tener el mismo orden. Preámbulo, Artículos, Disposiciones adicionales, Disposiciones Transitorias, Disposiciones Derogatorias y Disposiciones Finales. En caso de conflicto, la prioridad es en cuanto al contenido: Transitorio, Derogatorio, Final y Adicional. La prioridad no habla de disposiciones, puede haber texto transitorio en una disposición final.

Cada tipo de disposición puede tener uno o varios fines. Normalmente es texto añadido a la normativa.

\paragraph{Disposiciones Adicionales}
    Se usa para régimen jurídico adicional que no tenga un acomodo en el texto articulado. Estos tienen que tener conexión con el objeto de la norma y no pueden modificar derechos positivos (pues esto queda a las disposiciones finales).

    Excepciones a la aplicación de la norma o requisitos.

    Mandatos no dirigidos a la norma jurídica. Pueden ser mandatos de cote político o programático, no tienen un contenido claro ni legal, sino que pueden instar de forma genérica al gobierno o la administración a perseguir un fin.

    Preceptos residuales (cajón de sastre) que sea imposible encontrarles un lugar mejor dentro de la ley.

\paragraph{Disposiciones Transitorias}
    Se utilizan restrictivamente para condiciones de índole temporal establecido expresamente. Los preceptos que no incluyan el periodo de aplicación, o sólo se puedan aplicar una vez, no son transitorios, sino finales.
    
    Este puede ser la pervivencia de la regulación antigua para situaciones producidas con anterioridad a la norma. Normalmente los procedimientos iniciados anteriores a la norma.

    También puede ser la pervivencia de la regulación antigua para situaciones producidas posteriormente a la nueva ley. Esto se suele hacer para facilitar la transición de una norma a otra, especialmente cuando la nueva es mas restrictiva.

    Regulación expresa de situaciones producidas antes de las normas. Se trata de una regulación específica (ni la anterior ni la actual).

    Aplicación retroactiva de la nueva norma. Normalmente en situaciones favorables.

    Regulación \textit{ad hoc} para situaciones posteriores a la norma. Es una regulación específica que crea un régimen especial transitorio.

\paragraph{Disposiciones Derogatorias}
    Deberían ser precisas y expresas. Determinando con claridad las partes de las normas que se derogan y las que se mantienen vigentes (por separado). Para evitar que pervivan normas con el mismo ámbito de aplicación, es aconsejable que las disposiciones de la norma anterior no derogadas, se incorporen a la nueva como disposiciones adicionales o transitorias.
    
    Las derogaciones no pueden ser genéricas. Aunque esto es muy dificil en la realidad. La complegidad del ordenamiento jurídico no permite saber con exactitud las disposiciones que se derogan, especialmente si la nueva norma no sustituye en bloque a la vieja o afecta parcialmente a varias normas. Cuando se incluyen disposiciones derogatorias expresas, no suele recogerse qué preceptos quedan en vigor, sino que se añade la derogación genérica «quedan derogadas cuantas disposiciones de igual o inferior rango se opongan a lo dispuesto en esta norma».

    La ley de Procedimiento Administrativo obligaba a incorporar una tabla de derogaciones y vigencias, que raramente se ha cumplido.

\paragraph{Disposiciones Finales}
    Modifican excepcionalmente el derecho vigente que no constituye el objeto principal de la norma. También el de atribuir a ciertas disposiciones un rango diferente al de su norma. Es poeligroso aumentar el rango de una norma.

    También puede reglarse aquí un calendario de entrada en vigor de la norma, si es distinto del de los 20 días del Código Civil. Incluida una entrada en vigor escalonada indicando días y preceptos.

\subsection{Principios del derecho}

% \setcounter{section}{0}
\section{Esquema Nacional de Seguridad}

Nombre completo: Real Decreto 311-2022 de 3 de mayo por el que se regula el Esquema Nacional de Seguridad. Deroga al RD 3/2010 homónimo. Su última revisión es de 2024-11-06 

\localtableofcontents

\subsection{Forma}
Consta de diez capítulos y 44 artículos. Cuatro disposiciones adicionales, una transitoria, una derogatoria y tres finales